% !TEX root = ../main.tex
\chapter{Serialization and Safetensors}
\label{ch:serialization}

Saving and loading model state is one of those features that is
boring when it works and infuriating when it does not. This
chapter treats \Lucid's serialization infrastructure: the
on-disk format (safetensors as primary, legacy pickle for
backwards compatibility), the \code{state\_dict} contract,
the optimiser-state handling, the cross-version compatibility
strategy, the dtype-conversion-on-load policy, and the
production-deployment patterns. The chapter pairs with
\chref{ch:pretrained} (the hub itself) and uses the same
underlying format machinery.

\section{The state\_dict Contract}
\label{sec:serialization:state_dict}

\Lucid\ adopts the convention of representing a model's
state as a flat \code{dict[str, Tensor]} keyed by dot-separated
parameter paths. For a small model:

\begin{lstlisting}[style=lucid,language=LucidPython]
class Net(nn.Module):
    def __init__(self):
        super().__init__()
        self.fc1 = nn.Linear(10, 32)
        self.fc2 = nn.Linear(32, 10)

model = Net()
sd = model.state_dict()
# sd is:
# {
#     "fc1.weight": Tensor(shape=(32, 10), dtype=fp32),
#     "fc1.bias":   Tensor(shape=(32,),    dtype=fp32),
#     "fc2.weight": Tensor(shape=(10, 32), dtype=fp32),
#     "fc2.bias":   Tensor(shape=(10,),    dtype=fp32),
# }
\end{lstlisting}

The flat-dict format is the input to all serialization paths:
safetensors, pickle (legacy), checkpoint resume, hub upload.

\subsection{Buffers}
\label{ssec:serialization:buffers}

Some modules have non-parameter state: BatchNorm's running
mean/variance, position embeddings that are
buffer-not-parameter, etc. These appear in \code{state\_dict}
under the same flat-key convention but are flagged separately
in \code{model.named\_buffers()}.

\subsection{Loading}
\label{ssec:serialization:loading}

\begin{lstlisting}[style=lucid,language=LucidPython]
model.load_state_dict(sd, strict=True)
\end{lstlisting}

Behaviour:
\begin{itemize}
  \item \code{strict=True}: every key in the model must be in
    \code{sd}, and vice versa. Mismatch raises.
  \item \code{strict=False}: report mismatched keys, do not
    raise. Used for partial loads (e.g., loading a backbone
    into a model with new head).
  \item \code{assign=False}: copy the source tensor into the
    destination's storage (preserves the destination's
    autograd state).
  \item \code{assign=True}: replace the destination tensor
    object entirely.
\end{itemize}

\section{Safetensors as Primary Format}
\label{sec:serialization:safetensors_primary}

The primary on-disk format is safetensors \cite{safetensors}:

\begin{lstlisting}[style=lucid,language=LucidPython]
from lucid.serialization import save_safetensors, load_safetensors

save_safetensors(model.state_dict(), "model.safetensors")
sd = load_safetensors("model.safetensors")
model.load_state_dict(sd)
\end{lstlisting}

The format's structure:
\begin{enumerate}
  \item 8-byte little-endian length prefix giving the header
    length $N$.
  \item $N$-byte UTF-8 JSON header.
  \item Raw tensor data, concatenated in the order specified
    by the header.
\end{enumerate}

The JSON header:

\begin{lstlisting}[style=lucid,language=LucidPython]
{
    "fc1.weight": {
        "dtype": "F32",
        "shape": [32, 10],
        "data_offsets": [0, 1280]
    },
    "fc1.bias": {
        "dtype": "F32",
        "shape": [32],
        "data_offsets": [1280, 1408]
    },
    ...
    "__metadata__": {
        "lucid_version": "3.5.0",
        "model_id": "net",
        "created_utc": "2026-05-24T15:30:00Z"
    }
}
\end{lstlisting}

The \code{\_\_metadata\_\_} field is optional and tracks
provenance.

\subsection{Why safetensors over pickle}
\label{ssec:serialization:why_safetensors}

The pickle-based formats \cite{pickle_python} (legacy in many
ML projects) have three liabilities:

\begin{enumerate}
  \item \textbf{Code execution}: deserialising pickle runs
    arbitrary Python in the file. A malicious checkpoint can
    own your machine. Safetensors deserialises only tensors.
  \item \textbf{Memory-mapping}: pickle reads the entire file
    into memory. Safetensors can be mmap'd; tensors materialise
    lazily.
  \item \textbf{Cross-framework}: pickle is Python-specific.
    Safetensors is a published spec readable by any compatible
    library.
\end{enumerate}

\subsection{Load times}
\label{ssec:serialization:load_times}

\begin{table}[t]
\centering
\small
\begin{tabular}{lrrr}
\toprule
Checkpoint & Size (GB) & pickle (s) & safetensors (s) \\
\midrule
ResNet-50           & 0.1   & 0.3 & 0.05 (mmap)  \\
GPT-2 medium        & 1.4   & 4.2 & 0.10         \\
BERT-large          & 1.3   & 4.0 & 0.10         \\
LLaMA-7B            & 13.0  & 38.0 & 0.5         \\
\bottomrule
\end{tabular}
\caption[Load times: pickle vs safetensors.]{Checkpoint load
times on M3~Max NVMe SSD. Safetensors with mmap is essentially
instant; pickle scales with file size because it materialises
all bytes.}
\label{tab:serialization:load_times}
\end{table}

\section{Sharded Safetensors}
\label{sec:serialization:sharded}

For checkpoints above 4\,GB, the file is sharded:

\begin{itemize}
  \item \code{weights-00001-of-00007.safetensors}
  \item \code{weights-00002-of-00007.safetensors}
  \item \dots
  \item \code{weights.safetensors.index.json}: maps
    parameter name to file shard + offset.
\end{itemize}

Save:

\begin{lstlisting}[style=lucid,language=LucidPython]
from lucid.serialization import save_sharded_safetensors

save_sharded_safetensors(
    sd,
    "llama_7b",
    max_shard_size="2GB",
)
\end{lstlisting}

Load:

\begin{lstlisting}[style=lucid,language=LucidPython]
from lucid.serialization import load_sharded_safetensors

sd = load_sharded_safetensors("llama_7b")
\end{lstlisting}

The user does not see individual shards; the loader handles
the assembly.

\section{Optimiser State}
\label{sec:serialization:optimizer_state}

The optimiser also has state to save: per-parameter moments,
step counters, learning-rate schedules. \Lucid's optimisers
expose:

\begin{lstlisting}[style=lucid,language=LucidPython]
opt_sd = optimizer.state_dict()
# {
#     "state": {
#         0: {"step": 1500, "exp_avg": Tensor, "exp_avg_sq": Tensor},
#         1: {"step": 1500, ...},
#         ...
#     },
#     "param_groups": [{"lr": 1e-3, "betas": (0.9, 0.999), ...}],
# }
\end{lstlisting}

The integer keys in \code{state} refer to position in the
parameter list (in the order they were registered with the
optimiser). On load, the same positional mapping must hold ---
typically guaranteed if both save and load use the same
model architecture.

\subsection{Why optimiser state matters}
\label{ssec:serialization:optimizer_state_matters}

Resuming training without optimiser state means losing the
running moments; for Adam this typically causes a brief
training instability before the moments re-stabilise. For
critical resumes (long training runs), save both:

\begin{lstlisting}[style=lucid,language=LucidPython]
lucid.save({
    "model": model.state_dict(),
    "optimizer": optimizer.state_dict(),
    "scheduler": scheduler.state_dict(),
    "epoch": epoch,
    "global_step": global_step,
    "rng_state": lucid.random.get_rng_state(),
}, "checkpoint.pt")
\end{lstlisting}

The high-level \code{lucid.save} wraps safetensors for the
tensors and a small JSON sidecar for the non-tensor metadata.

\section{Checkpoint Format: lucid.save / lucid.load}
\label{sec:serialization:lucid_save}

The user-friendly API:

\begin{lstlisting}[style=lucid,language=LucidPython]
# Save a dict of mixed tensors and metadata
lucid.save({
    "model": model.state_dict(),
    "optimizer": opt.state_dict(),
    "global_step": 1500,
    "config": {"lr": 1e-3, "batch": 32},
}, "checkpoint.lpt")

# Load
ckpt = lucid.load("checkpoint.lpt")
model.load_state_dict(ckpt["model"])
opt.load_state_dict(ckpt["optimizer"])
\end{lstlisting}

The \code{.lpt} format:
\begin{itemize}
  \item Internally a tar archive.
  \item Contains \code{tensors.safetensors} (for all tensors).
  \item Contains \code{meta.json} (for non-tensor metadata).
\end{itemize}

The format is self-contained and reproducible.

\section{Cross-Version Compatibility}
\label{sec:serialization:cross_version}

A 3.5.0-saved checkpoint must load in 3.5.x and 3.6.x. The
policy:

\begin{itemize}
  \item \textbf{Forward compat (load older in newer)}:
    guaranteed within major versions. The loader handles
    renamed parameters via per-model translation tables.
  \item \textbf{Backward compat (load newer in older)}: not
    guaranteed. New tensors may simply be missing from the
    older model's state\_dict.
\end{itemize}

The \code{\_\_metadata\_\_} field tracks the version that
saved the checkpoint; the loader uses this to apply
version-specific transformations.

\subsection{Renaming policy}
\label{ssec:serialization:renaming_policy}

When a parameter is renamed in a new version, the model's
\code{load\_state\_dict} consults a renaming table:

\begin{lstlisting}[style=lucid,language=LucidPython]
class MyModel(nn.Module):
    state_dict_rename_table = {
        # Pre-3.5 used "linear" instead of "fc"
        "linear.weight": "fc.weight",
        "linear.bias":   "fc.bias",
    }
\end{lstlisting}

The loader applies the table during \code{load\_state\_dict}.
The table is maintained alongside the model code; deletion
requires a major-version bump.

\section{Dtype Conversion on Load}
\label{sec:serialization:dtype_conversion}

A checkpoint saved at FP32 should load into an FP16 model
without manual cast. \Lucid's policy:

\begin{itemize}
  \item By default, the loader casts to the destination's
    dtype.
  \item \code{load\_state\_dict(sd, dtype=lucid.float32)}
    forces all loaded tensors to FP32.
  \item \code{load\_state\_dict(sd, dtype="match")} (default)
    matches the destination parameter's dtype.
\end{itemize}

Implementation detail: the cast happens on the destination
device, not on disk. The safetensors mmap brings bytes into
unified memory; the cast is a fast on-GPU operation.

\section{Tensor Cloning vs Sharing}
\label{sec:serialization:clone_vs_share}

By default, \code{load\_safetensors} returns mmap'd tensors:
the tensor's storage is the file's mmap region. This is fast
but creates a subtle issue: if the user modifies the tensor in
place, the modification is visible to any other reader of the
same file (and crashes if the file is read-only).

The policy:
\begin{itemize}
  \item For \code{model.load\_state\_dict}: tensors are cloned
    on copy (no shared mmap).
  \item For direct \code{load\_safetensors} return: tensors
    are mmap-backed. The user is responsible for cloning
    before mutation.
\end{itemize}

The \code{copy=True} flag forces eager cloning:

\begin{lstlisting}[style=lucid,language=LucidPython]
sd = load_safetensors("model.safetensors", copy=True)
# All tensors are owned-storage, safe to mutate
\end{lstlisting}

\section{Atomic Writes}
\label{sec:serialization:atomic_writes}

The \code{lucid.save} writes atomically: it writes to a
temporary file in the same directory, then renames over the
target. This ensures that a crash during the write leaves the
old file intact, not a corrupted half-write.

\begin{lstlisting}[style=lucid,language=LucidPython]
# Internally:
tmp = target + ".tmp." + str(os.getpid())
with open(tmp, "wb") as f:
    write_safetensors(f, state_dict)
    f.flush()
    os.fsync(f.fileno())
os.replace(tmp, target)
\end{lstlisting}

The same pattern protects checkpoint files written during
training: a power failure during save will not corrupt the
checkpoint that was already on disk.

\section{Checkpoint Compression}
\label{sec:serialization:compression}

For deployment-size optimisation, \Lucid\ supports an optional
quantisation-then-save path:

\begin{lstlisting}[style=lucid,language=LucidPython]
lucid.save(model.state_dict(), "model.lpt", dtype=lucid.float16)
# Casts all tensors to FP16 before saving; halves file size.
\end{lstlisting}

For more aggressive quantisation (INT8 weights), see
\code{lucid.serialization.quant}. Note: quantised checkpoints
must be loaded into a quantised model (mismatched dtype on
load is not a free conversion in general).

\section{ONNX Export}
\label{sec:serialization:onnx}

For cross-framework deployment, \Lucid\ provides ONNX
\cite{onnx_2017} export:

\begin{lstlisting}[style=lucid,language=LucidPython]
from lucid.export.onnx import export_onnx

export_onnx(
    model,
    args=(lucid.randn(1, 3, 224, 224),),
    output_path="model.onnx",
    opset_version=18,
)
\end{lstlisting}

The exporter:
\begin{enumerate}
  \item Traces the model with the provided example inputs.
  \item Maps each op in the trace to its ONNX equivalent.
  \item Writes the result.
\end{enumerate}

Unsupported ops raise with a clear message (the user can
either replace the op or open an issue requesting ONNX
support).

\section{Core ML Export}
\label{sec:serialization:coreml}

For Apple-platform deployment (iOS / macOS Vision API),
\Lucid\ provides Core ML export:

\begin{lstlisting}[style=lucid,language=LucidPython]
from lucid.export.coreml import export_coreml

export_coreml(
    model,
    inputs=[lucid.export.coreml.TensorSpec(shape=(1, 3, 224, 224))],
    output_path="model.mlmodel",
)
\end{lstlisting}

The Core ML model can be loaded on-device for inference using
the Vision framework. The export pipeline shares the
ONNX exporter's tracing infrastructure.

\section{Format Conversion Tools}
\label{sec:serialization:format_conversion}

For users with checkpoints in other formats:

\begin{lstlisting}[style=lucid,language=LucidPython]
$ lucid convert --from pytorch_pickle model.pt --to safetensors model.safetensors
$ lucid convert --from hf model_name --to safetensors model.safetensors
\end{lstlisting}

The \code{convert} CLI handles common formats. For exotic
ones, the user writes a small Python script reading the source
and calling \code{save\_safetensors}.

\section{Summary and Forward References}
\label{sec:serialization:summary}

\Lucid's serialization stack: safetensors as the primary
on-disk format (fast, safe, mmap-friendly), \code{state\_dict}
as the in-memory contract, atomic writes for crash safety,
cross-version renaming for forward compatibility, dtype
conversion on load, ONNX and Core ML exports for deployment.

The next chapter, \chref{ch:data_pipeline}, treats the data
pipeline. \chref{ch:error_handling} treats error messaging.
\chref{ch:reproducibility} treats reproducibility.

\begin{lucidnote}
For the user: \code{lucid.save} / \code{lucid.load} is the
canonical checkpoint API; \code{save\_safetensors} is for
weights-only saves. Always save optimiser state for long
training runs. For deployment, prefer safetensors over the
legacy pickle path.
\end{lucidnote}
